\documentclass[11pt, a4paper]{article}
\usepackage[a4paper, margin=2.5cm]{geometry}
\usepackage{fontspec}
\usepackage[english, bidi=basic, provide=*]{babel}
\babelprovide[import, onchar=ids fonts]{english}
\babelfont{rm}{Noto Sans}

\usepackage{amsmath, amssymb}
\usepackage{booktabs}
\usepackage{array}
\usepackage{xcolor}
\usepackage{hyperref}

\definecolor{primary}{RGB}{30, 64, 175}
\definecolor{secondary}{RGB}{100, 116, 139}
\definecolor{lightbg}{RGB}{248, 250, 252}

\hypersetup{
    colorlinks=true,
    linkcolor=primary,
    urlcolor=primary
}

\setlength{\parindent}{0pt}
\setlength{\parskip}{8pt}

\title{\textbf{\color{primary} Comprehensive Notes: Graphing Trigonometric Functions}}
\author{Trigonometry Study Guide}
\date{\today}

\begin{document}

\maketitle

\section{Standard Form of Trigonometric Functions}

Any transformational sine or cosine function can be written in the general form:
\begin{equation*}
y = A \sin(B(x - C)) + D \quad \text{or} \quad y = A \cos(B(x - C)) + D
\end{equation*}

Where each parameter dictates a specific geometric transformation on the parent graph:

\begin{itemize}
    \item \textbf{Amplitude ($|A|$):} The vertical stretch or compression. Represents half the distance between the maximum and minimum values. If $A < 0$, the graph is reflected across the midline.
    \item \textbf{Period ($T$):} The distance required for the function to complete one full cycle.
    \begin{equation*}
    \text{Period } (T) = \frac{2\pi}{|B|}
    \end{equation*}
    \item \textbf{Phase Shift ($C$):} The horizontal displacement of the function.
    \begin{itemize}
        \item If $C > 0$, the graph shifts to the \textbf{right} by $C$ units.
        \item If $C < 0$, the graph shifts to the \textbf{left} by $|C|$ units.
    \end{itemize}
    \item \textbf{Vertical Shift / Midline ($D$):} The horizontal line $y = D$ around which the wave oscillates.
\end{itemize}

---

\section{Key Characteristics of Parent Functions}

\subsection{Sine and Cosine}

\begin{table}[h]
\centering
\begin{tabular}{l p{5cm} p{5cm}}
\toprule
\textbf{Property} & \textbf{Sine Function } $y = \sin(x)$ & \textbf{Cosine Function } $y = \cos(x)$ \\
\midrule
Domain & $(-\infty, \infty)$ & $(-\infty, \infty)$ \\
Range & $[-1, 1]$ & $[-1, 1]$ \\
Period & $2\pi$ & $2\pi$ \\
Amplitude & $1$ & $1$ \\
Symmetry & Odd (Symmetric about origin) & Even (Symmetric about y-axis) \\
Key Points (1 Cycle) & $(0,0), (\frac{\pi}{2},1), (\pi,0), (\frac{3\pi}{2},-1), (2\pi,0)$ & $(0,1), (\frac{\pi}{2},0), (\pi,-1), (\frac{3\pi}{2},0), (2\pi,1)$ \\
\bottomrule
\end{tabular}
\caption{Comparison of parent Sine and Cosine functions.}
\end{table}

\subsection{Tangent and Reciprocal Functions}

\begin{itemize}
    \item \textbf{Tangent ($y = \tan(x)$):} 
    \begin{itemize}
        \item Period: $T = \frac{\pi}{|B|}$
        \item Vertical Asymptotes: $x = \frac{\pi}{2} + k\pi$ for any integer $k$.
        \item Domain: $x \neq \frac{\pi}{2} + k\pi$; Range: $(-\infty, \infty)$.
    \end{itemize}
    \item \textbf{Cosecant ($y = \csc(x) = \frac{1}{\sin(x)}$):}
    \begin{itemize}
        \item Period: $2\pi$; Vertical Asymptotes at $x = k\pi$.
        \item Local minima/maxima touch the peak and trough points of $y = \sin(x)$.
    \end{itemize}
    \item \textbf{Secant ($y = \sec(x) = \frac{1}{\cos(x)}$):}
    \begin{itemize}
        \item Period: $2\pi$; Vertical Asymptotes at $x = \frac{\pi}{2} + k\pi$.
    \end{itemize}
\end{itemize}

---

\section{Step-by-Step Graphing Strategy}

To accurately sketch a function such as $y = A \sin(B(x - C)) + D$:

\begin{enumerate}
    \item \textbf{Identify Core Parameters:}
    Extract $A$, $B$, $C$, and $D$ from the given equation (factor out $B$ inside the argument if necessary).
    
    \item \textbf{Determine Key Attributes:}
    \begin{itemize}
        \item Midline: $y = D$
        \item Maximum Value: $y_{\text{max}} = D + |A|$
        \item Minimum Value: $y_{\text{min}} = D - |A|$
        \item Period: $T = \frac{2\pi}{|B|}$
        \item Increment (Quarter-Period): $\Delta x = \frac{T}{4}$
    \end{itemize}
    
    \item \textbf{Calculate Key X-Values:}
    Start at the phase shift $x_0 = C$, and consecutively add the quarter-period increment $\Delta x$:
    \begin{equation*}
    x_0 = C, \quad x_1 = C + \Delta x, \quad x_2 = C + 2\Delta x, \quad x_3 = C + 3\Delta x, \quad x_4 = C + 4\Delta x
    \end{equation*}

    \item \textbf{Evaluate Corresponding Y-Values:}
    \begin{itemize}
        \item \textbf{Sine Pattern:} Midline $\rightarrow$ Max $\rightarrow$ Midline $\rightarrow$ Min $\rightarrow$ Midline
        \item \textbf{Cosine Pattern:} Max $\rightarrow$ Midline $\rightarrow$ Min $\rightarrow$ Midline $\rightarrow$ Max
        \item \textit{(Note: If $A < 0$, invert the Max and Min points.)}
    \end{itemize}

    \item \textbf{Plot and Smooth Curve:}
    Plot the five primary points, draw vertical asymptotes if applicable, and connect points with a continuous smooth wave.
\end{enumerate}

---

\section{Worked Example}

\textbf{Problem:} Graph the function $y = -3 \cos\left(2x - \frac{\pi}{2}\right) + 1$ over one full period.

\subsection*{Solution Walkthrough}

\textbf{Step 1: Rewrite in factored standard form}
\begin{equation*}
y = -3 \cos\left(2\left(x - \frac{\pi}{4}\right)\right) + 1
\end{equation*}

\textbf{Step 2: Identify parameters}
\begin{itemize}
    \item $A = -3$ (Amplitude = $3$, reflected over midline)
    \item $B = 2 \implies \text{Period } T = \frac{2\pi}{2} = \pi$
    \item $C = \frac{\pi}{4}$ (Phase shift right by $\frac{\pi}{4}$)
    \item $D = 1$ (Midline at $y = 1$)
\end{itemize}

\textbf{Step 3: Calculate increments and key points}
\begin{itemize}
    \item Quarter-Period Increment: $\Delta x = \frac{T}{4} = \frac{\pi}{4}$
    \item $x_0 = \frac{\pi}{4}$
    \item $x_1 = \frac{\pi}{4} + \frac{\pi}{4} = \frac{\pi}{2}$
    \item $x_2 = \frac{\pi}{2} + \frac{\pi}{4} = \frac{3\pi}{4}$
    \item $x_3 = \frac{3\pi}{4} + \frac{\pi}{4} = \pi$
    \item $x_4 = \pi + \frac{\pi}{4} = \frac{5\pi}{4}$
\end{itemize}

\textbf{Step 4: Determine point coordinates}

Since $A = -3 < 0$, the inverted cosine pattern is: **Min $\rightarrow$ Midline $\rightarrow$ Max $\rightarrow$ Midline $\rightarrow$ Min**.

\begin{table}[h]
\centering
\begin{tabular}{c c p{6cm}}
\toprule
$x$ & $y$ & \textbf{Key Feature} \\
\midrule
$\frac{\pi}{4}$ & $1 - 3 = -2$ & Minimum point \\
$\frac{\pi}{2}$ & $1$ & Midline intercept \\
$\frac{3\pi}{4}$ & $1 + 3 = 4$ & Maximum point \\
$\pi$ & $1$ & Midline intercept \\
$\frac{5\pi}{4}$ & $1 - 3 = -2$ & Minimum point \\
\bottomrule
\end{tabular}
\caption{Key coordinates for $y = -3 \cos\left(2\left(x - \frac{\pi}{4}\right)\right) + 1$ over one period.}
\end{table}

\end{document}